\begin{proposition}[4.9.4]
\label{IV.2.4.9.4-fr}
Con las notaciones de \textup{(4.9.1)}, supongamos que $K$ sea una extensión separable
de $K'$~; para que $K'$ sea un cuerpo de definición de $\overline T$, es
necesario y suficiente que $K'$ sea un cuerpo de definición del subpreesquema
reducido de $X_{(K)}$ cuyo espacio subyacente es $\overline T$.
\end{proposition}

\begin{proof}
La condición es evidentemente suficiente en virtud de (4.9.1)~; para ver que
es necesaria, observemos que, si $\overline T=g^{-1}(T')$, donde $T'$ es una
parte cerrada de $X_{(K')}$, y si $Z'$ es el subpreesquema reducido de
$X_{(K')}$ cuyo espacio subyacente es $T'$, entonces
$\overline Z=Z'\otimes_{K'}K=\operatorname{Spec}(K)\times_{K'}Z'$ es reducido
en virtud del criterio (4.6.1, \emph{e)}) y de la hipótesis sobre $K$, y tiene como
espacio subyacente $\overline T$.
\end{proof}

\begin{corollary}[4.9.5]
\label{IV.2.4.9.5-fr}
Supongamos que se satisface una de las hipótesis siguientes:
\begin{enumerate}
\item[\textup{a)}] $k$ es de característica $0$.
\item[\textup{b)}] $K$ es una extensión algebraica de $k$ y $X$ es de tipo finito sobre
$k$.
\end{enumerate}
Entonces toda parte cerrada $\overline T$ de $X_{(K)}$ posee un menor cuerpo
de definición, que es una extensión separable de $k$ (y finita sobre
$k$ bajo la hipótesis \textup{b))}.
\end{corollary}

\begin{proof}
Bajo la hipótesis \emph{a)}, $K$ es separable sobre cada uno de sus subcuerpos, por lo
que de (4.9.4) se deduce que los cuerpos de definición de $\overline T$ son los
mismos que los del subpreesquema reducido de $X_{(K)}$ cuyo espacio subyacente
es $\overline T$. El corolario se deduce entonces de (4.8.11). Bajo la
hipótesis \emph{b)}, se sabe que $K$ es una extensión radicial de una extensión
separable $K_1$ de $k$. Entonces, para toda subextensión $K'$ de $K$, $K'$ es
radicial sobre $K'\cap K_1$ y, por tanto, en virtud de (4.9.2), $K'$ es un
cuerpo de definición de $\overline T$ si y solo si $K'\cap K_1$ lo es. Esto
nos reduce al caso en que $K=K_1$, y por tanto al caso en que $K'$ es una
extensión \emph{algebraica separable} de $k$, y el razonamiento se termina como en el
caso \emph{a)}.
\end{proof}

\begin{proposition}[4.9.6]
\label{IV.2.4.9.6-fr}
Sean $k$ un cuerpo, $X$ un $k$-preesquema, $K$ una extensión de $k$,
$\overline U$ una parte a la vez abierta y cerrada de $X_{(K)}$, y
$\overline V$ su complementaria. Sea $K'$ una subextensión de $K$~; entonces
las condiciones siguientes son equivalentes:
\begin{enumerate}
\item[\textup{a)}] $K'$ es un cuerpo de definición de la parte $\overline U$ de
$X_{(K)}$.
\item[\textup{a')}] $K'$ es un cuerpo de definición del subpreesquema cerrado de
$X_{(K)}$ inducido sobre el abierto $\overline U$.
\item[\textup{b)}] $K'$ es un cuerpo de definición de la parte $\overline V$ de
$X_{(K)}$.
\item[\textup{b')}] $K'$ es un cuerpo de definición del subpreesquema cerrado de
$X_{(K)}$ inducido sobre el abierto $\overline V$.
\end{enumerate}
\end{proposition}

\begin{proof}
Es claro que \emph{a')} implica \emph{a)}, y \emph{a)} implica \emph{b')} en virtud de (4.9.1). Como
$\overline U$ y $\overline V$ desempeñan papeles simétricos, \emph{b')} implica \emph{b)} y
\emph{b)} implica \emph{a')}, lo que termina la demostración.
\end{proof}

\begin{corollary}[4.9.7]
\label{IV.2.4.9.7-fr}
Bajo las hipótesis de \textup{(4.9.6)}, existe un menor cuerpo de definición $K'$ para
la parte $\overline U$ de $X_{(K)}$, y $K'$ es también el menor cuerpo de
definición de la parte $\overline V$ de $X_{(K)}$ y de los
subpreesquemas cerrados de $X_{(K)}$ inducidos sobre los abiertos $\overline U$
y $\overline V$. Si además $X$ es de tipo finito sobre $k$, $K'$ es una
extensión finita separable de $k$.
\end{corollary}

\begin{proof}
Teniendo en cuenta (4.9.6) y (4.8.11), solo queda demostrar la última
afirmación. En virtud de (4.8.15, (ii)), es posible limitarse al caso en que $K$
es algebraicamente cerrado. Sea entonces $\overline k$ la clausura algebraica
de $k$ contenida en $K$, y denotemos por $W_i$ las componentes

\oldpage[IV-2]{86}

conexas de $X_{(\overline k)}$ ($1\leq i\leq m$)~; si
$p:X_{(K)}\to X_{(\overline k)}$ es la proyección canónica, se sabe que los
$p^{-1}(W_i)$ son las componentes conexas de $X_{(K)}$ (4.4.6), por lo que
$\overline U$ es una unión de cierto número de estas componentes y, en
consecuencia, está definida sobre $\overline k$~; la conclusión se deduce
entonces de (4.9.5, \emph{b)}).
\end{proof}
